\section{Positive boundary mass and winding polygons}\label{sec:cylinder}

We prove two estimates with constants uniform in the circumference: a
boundary deficit of order \(N^{-1/4}\), and a mass of order \(N^{-1}\)
for one winding polygon whose highest row is fixed.  The first step is
to prove that the boundary sums over all depths are finite.  We then
express the mass of one polygon as a polynomial in a tile weight, constrain that
polynomial by fusion, and recover both estimates from a strip integral.

Retain \(l=\pi/8\) and \(x=(2\cos l)^{-1}\), and put \(c=3/8\).
Orient the unit triangular tiling with horizontal sides.  A boundary
port is the midpoint of a triangle side.  A port-to-port path visits
each triangle at most once and joins two distinct side midpoints through
its centre; a path visiting \(m\) centres has weight \(x^m\).
For an integer \(g>0\), let \(K_\infty(g)\) be the total weight of
paths in the lower half-plane from horizontal boundary port \(0\) to
port \(g\).  In the quotient by horizontal translation by an integer
\(N\ge1\), let \(K_N(g)\), \(1\le g<N\), be the corresponding sum over
paths that are self-avoiding in the quotient and whose chosen lift has
endpoints \(0,g\).  In particular, \(K_N(g)\le K_\infty(g)\).
All domain boundaries and row cuts in this section follow triangle
sides.

Group the triangles into rows of \(\pi/3\)-rhombi, with sides parallel
to \(0\) and \(\pi/3\); the row index increases by one under translation
by \(e^{i\pi/3}\).  Let \(P_N\) be the sum of \(x^{|\gamma|}\) over
simple noncontractible polygons in the infinite quotient cylinder
whose highest occupied row is a specified row.  Here \(|\gamma|\) is
the number of visited centres, equivalently the number of polygon
edges.  Each actual quotient polygon is counted once without
orientation.  Thus \(P_N\) is a bare sum of single-polygon weights; it
will be the derivative at winding fugacity zero, with no factor two.

\begin{proposition}[Cylinder deficit and winding mass]\label{prop:cylinder}
In these port and highest-row conventions,
\[
 D_N:=1-2\cos(3l)\sum_{g=1}^{N-1}K_N(g)\asymp N^{-1/4},
 \qquad
 P_N\asymp N^{-1}.
\]
The comparison constants can be chosen independently of the integer
circumference \(N\ge1\).
\end{proposition}

\subsection{Finite boundary flux and half-plane summability}

Consider a convex finite domain tiled by triangles, and fix an entrance
port \(a\).  Give a path of total signed turn \(W\) the complex weight
\(e^{icW}x^m\).  The sum of these weights over paths from \(a\) to all
exit ports is \(1\).  We recall the local cancellation, including its
phases, because the sign changes when the source lies in an interior
slit later in the paper.

Introduce one unit of inward flux at \(a\), and stop a partial path
whenever it reaches a side midpoint, recording whether it points into
a triangle or out of one.  At a fresh triangle the two extensions turn
by \(\pm\pi/3\).  Their total multiplier is
\[
 2x\cos(c\pi/3)=1.
\]
At a previously visited triangle, pair the two stopped paths obtained
by reversing the cycle segment.  To check their phases, suppose that
the earlier arrival-to-exit turn was \(+\pi/3\).  The closing turn of
the cycle is then \(-\pi/3\).  The prefix from the boundary source
arrives from the exterior of that cycle, so the cycle has rotation
\(-2\pi\): in the counterclockwise alternative the reflex interior
sector would contain the arriving prefix.  Including the earlier turn
gives a net increment \(-4\pi/3\), which becomes \(+4\pi/3\) on
reversal.  The pair cancels because
\[
 2\cos(4\pi c/3)=0.
\]
The reflected case has the same cancellation.  Thus all interior
stops cancel and the input flux equals the boundary output.

For an exit port \(b\) following \(a\) counterclockwise, let
\(\phi_a,\phi_b\) be the lifted counterclockwise boundary tangents.
Closing the path along the boundary gives
\[
 W=\phi_b-\phi_a-\pi.
\]
Convexity puts \(W\) in \([-\pi,\pi]\), so every real output weight is
positive.  Exhaust the lower half-plane by increasing triangles
symmetric about \(a\).  The two directions along their horizontal
base have phases \(e^{\pm ic\pi}\), both with real part
\(\cos(3l)\).  Monotone convergence for these base paths, with the
other positive exits discarded, gives
\begin{equation}\label{eq:cyl-half-plane}
 2\cos(3l)\sum_{g>0}K_\infty(g)\le1.
\end{equation}
This finite bound is the input for the annular cap construction.  The
equality and the endpoint tail estimate will follow after the deficit
has been evaluated.

\subsection{Annular diagram rows and the zero-fugacity cap}

We first distinguish the two row operators that enter the argument.
The spin row \(T^\zeta(u;\mathbf b)\) of Section~\ref{sec:local} acts on
\(V^{\otimes N}\), with a twist \(\zeta^{s_a}\) in the auxiliary trace.
The diagram row defined below acts instead on annular arch patterns.
The two descriptions give the same scalars when the open diagram ports
are closed with the specified caps and the orientations are summed;
we do not identify their state spaces.

A tile has two incoming legs, auxiliary from the right and quantum
from below, and the opposite outgoing legs.  Its scalar weights are
the functions \(a,b,h,d,e\) of Section~\ref{sec:local}.  An empty tile
has weight \(1\).  A single connection from one input to the other
stream's output has weight \(a\), a single straight connection has
weight \(h\), and a cup or cap joining the two inputs or the two
outputs has weight \(b\).  When all four legs are occupied, the
double pairing of \(a\)-type has weight \(d\), and the other double
pairing has weight \(e\).  Contractible closed loops have weight zero,
and simple winding loops have weight \(n\).  At the physical argument
\(t=l\),
\[
 (a,b,h,d,e)(l)=(x,x^2,x^2,x^2,0).
\]
These are exactly the weights of dual self-avoiding paths in the two
triangles of a split rhombus.

Here is the orientation comparison with the spin matrix.  Choose
\(\zeta\ne0\) with \(n=\zeta+\zeta^{-1}\) and orient each occupied
arc.  The spins \(+1,-1\) mean respectively along and against the
positive stream direction.  Straighten the positive stream tangents
at the tile ports to parallel forward directions, with the auxiliary
stream to the right of the quantum stream at entry and to its left
at exit.  A through connection then has net turn zero.  A clockwise
cup or cap has phase \(p=e^{-i\pi/4}\), and a counterclockwise one has
phase \(p^{-1}\), exactly as in the terms involving \(C\) and
\(p^{-m}\) in \(R(t)\).  These spin phases use one quarter of the turn
angle, separately from the boundary flux phase \(e^{icW}\) above.
The changes of tangent used in straightening cancel at each join, so
a closed curve of signed rotation \(\Theta\) has phase
\(e^{i\Theta/4}\).  Equivalently, the phase counts signed cups and
caps and is unchanged by isotopy fixing the port tangents.

A contractible simple loop has rotation \(\pm2\pi\); its two
orientations therefore contribute \(i\) and \(-i\), with sum zero.
A simple winding loop has rotation zero in the flat cylinder frame
and seam flux \(\pm1\); its orientation sum is
\(\zeta+\zeta^{-1}=n\).  This remains true for a diagram closed by
arches in earlier rows.  For each such arch one chooses an actual
representative below the row, with the prescribed port tangents and
homotopy, and includes that representative's own turn and seam
factors.  Representatives above a row are handled in the same way.

We also need the interchange and inversion identities with open
connections.  For interchange, place streams \(1,2,3\) from right
to left at entry, all pointing forward, and in reverse order at exit.
For fixed occupied boundary slots in a disk, the
oriented tensors of distinct noncrossing matchings are linearly
independent.  Indeed, after straightening the port tangents, the
nonzero arc phase depends only on the oriented connectivity.  For a
chosen matching, assign inward spins to its left endpoints and outward
spins to its right endpoints in a linearized boundary order.  The
result is a Dyck word.  At every gap, a matching compatible with this
assignment has at least the Dyck height many open arcs; equality at
every gap singles out the chosen matching.  Ordering matchings by
these height tests gives a triangular test matrix with nonzero
diagonal.  The spin identities of Section~\ref{sec:local} therefore
imply the corresponding identities of open diagrams.

Sliding a crossing of two auxiliary streams around the \(N\) sites
now proves commutation of periodic diagram rows: insert the crossing
and its inverse, use interchange at each site, and cancel them after
the trace.  Any loops formed during this operation are annular loops
with the prescribed weights.  The proof applies at generic regular
parameters and extends as a meromorphic identity.

Write \(\mathsf T_n(t)\) for the homogeneous diagram row with argument
\(t\) and winding weight \(n\).  A state specifies, at the bottom
boundary circle, which of the \(N\) slots are vacant and how the
occupied slots are joined by noncrossing arches below the circle,
including their homotopy classes.  No strand runs to infinity.
A chosen lift of each arch has endpoint span of absolute value less
than \(N\), since an arch disjoint from its translates cannot have
interlacing lifted endpoints.  There are consequently only finitely
many states.  To act by \(\mathsf T_n(t)\), attach a row above the
input arches, evaluate closed loops, and read the output arch pattern.
The preceding diagram interchange gives, for the same winding weight,
\[
 \mathsf T_n(u)\mathsf T_n(v)=\mathsf T_n(v)\mathsf T_n(u).
\]

Let \(|0\rangle\) be the all-vacant state and let \(\langle0|\) select
its component.  Set \(T(t):=\mathsf T_0(t)\), retaining this shorthand
for the zero-fugacity diagram row.  Every nonempty configuration with
empty output contains a closed loop, so
\[
 \langle0|T(t)=\langle0|.
\]
We next prove, for each fixed \(N\),
\begin{equation}\label{eq:cyl-cap-limit}
 T(l)^H\longrightarrow|\Psi\rangle\langle0|
 \qquad(H\to\infty),\qquad \langle0|\Psi\rangle=1.
\end{equation}

\paragraph{Empty input.}
The column \(T(l)^H|0\rangle\) is the positive sum of physical arch
systems of depth at most \(H\).  For a fixed output pattern, lift each
component path to the half-plane.  Its endpoint separation has
absolute value less than \(N\), and its total weight is bounded by
the corresponding \(K_\infty(g)\).  Ignoring mutual avoidance bounds
the arch-system sum by products of these finite masses.  Only
finitely many endpoint and homotopy choices occur at fixed \(N\).
Thus each component increases to a finite limit.  These limits define
the finite vector \(\Psi\); its empty component is \(1\).

\paragraph{Nonempty input.}
Fix a nonempty annular pattern \(\alpha\).  It has a physical
realization of positive weight in finitely many rows below the cut.
To construct one, repeatedly choose a pair whose bounded disk side
contains no endpoints still to be joined.  Cap that pair across its
auxiliary interval, using one \(a\), one \(b\), and the intervening
\(h\)'s in bottom-to-top notation, while carrying all other
occupations straight.  This also permits an interval crossing the
seam.  The construction terminates, and all its weights at \(l\)
are positive.

Attach this one fixed realization beneath a diagram counted by the
column \(T(l)^H|\alpha\rangle\).  The union determines the original
bulk diagram, so this attachment is injective.  At zero fugacity no
closed component survives.  Hence every segment belongs to an output
arch, and at least one such arch reaches the attached realization,
at depth at least \(H-O_\alpha(1)\).  The output weight multiplied by
the fixed positive cap weight is therefore bounded by the tail of the
convergent sum of completed arch systems.  That tail tends to zero.  Every
nonempty input column thus tends to zero, proving
\eqref{eq:cyl-cap-limit} in the finite state space.

Commutation and the left empty identity now give the common right
fixed vector:
\[
 T(t)|\Psi\rangle\langle0|
   =\lim_{H\to\infty}T(l)^H T(t)
   =|\Psi\rangle\langle0|T(t)
   =|\Psi\rangle\langle0|,
 \qquad T(t)\Psi=\Psi.
\]
The limit \eqref{eq:cyl-cap-limit} also says that \(1\) is a simple
eigenvalue of \(T(l)\), with all remaining eigenvalues strictly inside
the unit disk; the principal nonempty block has spectral radius less
than \(1\).  These are fixed-\(N\) conclusions.  They give a locally
analytic eigenpair under a small change of winding fugacity at each
fixed \(N\), as used for the higher derivatives in
Section~\ref{sec:projection}, but give no uniform spectral gap or
uniform analyticity radius in \(N\).

\subsection{The one-winding derivative, fusion, and its polynomial}

Keep the cap \(\Psi\) at zero fugacity and define
\begin{equation}\label{eq:cyl-bare-J}
 J_N(t):=
 \left\langle0\left|
       \left.\partial_n\mathsf T_n(t)\right|_{n=0}
 \right|\Psi\right\rangle .
\end{equation}
The derivative selects exactly one winding loop closed by the new
row: \(\partial_n n|_{0}=1\), while higher powers of \(n\) vanish.
The cap below contains ordinary open arches and no closed loops.
At \(t=l\) the new row is the unique highest occupied row of the
resulting polygon.  Thus each actual unoriented quotient polygon is
counted once, without a marked vertex or a division by horizontal
translations, and
\[
 J_N(l)=P_N.
\]
The fugacity-\(2\) gas used later instead assigns \(2x^{|\gamma|}\)
to each polygon.

\paragraph{Fusion after closing the caps.}
For any fixed annular bottom cap \(\alpha\), the local fusion identity
implies a scalar relation before differentiation:
\begin{equation}\label{eq:cyl-capped-fusion}
 \left\langle0\left|
   \mathsf T_n(t)\mathsf T_n(t+2l)-\mathsf T_n(t+l)
 \right|\alpha\right\rangle
 =t^N R_\alpha(t,\zeta),
 \qquad n=\zeta+\zeta^{-1},
\end{equation}
where \(R_\alpha\) is analytic near \(t=0,\zeta=i\).
To see the order \(N\), label the two auxiliary spaces \(1,2\).
The local product
\(R_{23}(t)R_{13}(t+2l)\) preserves
\(E_2V\otimes V_3\) and acts there as the single tile \(R(t+l)\).
Moreover, \(E_2\) intertwines \(\zeta^{s_1+s_2}\) with
\(\zeta^s\), so the trace on that subspace is the fused row.  On the
quotient auxiliary space, each local product is \(O(t)\): at \(t=0\)
it is \(P_{23}R_{13}(2l)\), whose range already lies in
\(E_2V\otimes V_3\).  Its product around \(N\) sites is therefore
\(O(t^N)\).  Closing with the chosen oriented representative of
\(\alpha\), including its turn and seam factors, converts this spin
trace relation to \eqref{eq:cyl-capped-fusion}.

There are finitely many cap patterns, each with a finite coefficient
in \(\Psi\), so \eqref{eq:cyl-capped-fusion} also holds with
\(\alpha\) replaced by \(\Psi\).  Near \(\zeta=i\), the quantity
\(n=\zeta+\zeta^{-1}\) is a local analytic coordinate because
\(dn/d\zeta=2\) there.  Differentiate that scalar relation at \(n=0\),
with the cap coefficients fixed.  The identities
\(T(t)\Psi=\Psi\) and \(\langle0|T(t)=\langle0|\) reduce the two
derivatives of the product to \(J_N(t)\) and \(J_N(t+2l)\).  Hence
\begin{equation}\label{eq:cyl-differentiated-fusion}
 J_N(t)+J_N(t+2l)-J_N(t+l)=O(t^N)
 \qquad(t\to0).
\end{equation}
No derivative of \(\Psi\) is taken in this argument.

\paragraph{Dependence on the straight weight.}
At empty output, a full horizontal auxiliary loop contributes
\(h(t)^N\) once.  Otherwise the occupied auxiliary intervals have
positive integer lengths \(s_1,\ldots,s_r\) and alternate with gaps
of positive lengths \(g_1,\ldots,g_r\).  Their bottom endpoints are
paired by arches in \(\Psi\).  We claim that each such arch crosses
one consecutive gap.

Indeed, in the lift an arch cuts off a bounded disk side.  It encloses
an even number of other endpoints, so its two endpoints are of
opposite types in the alternating list of interval starts and ends.
The incident row intervals therefore both point into that disk side
or both point out.  The first case closes a component inside the
bounded lifted disk.  In the second case, any enclosed endpoints are
paired internally both by arches and by row intervals, again making
a closed component inside that disk.  Such components are
contractible on the cylinder and have zero weight.  The disk side must
therefore contain no other endpoint, which is precisely a pairing
across a consecutive gap.

Put \(C_0=2\cos(3l)\).  The local weight formulas give, for each
occupied interval,
\[
 a(t)b(t)h(t)^{s_i-1}=C_0(1-h(t))h(t)^{s_i}.
\]
A labelled cyclic arrangement \(\mathcal A\) records its actual
intervals on the \(N\)-site ring and the resulting arches across the gaps;
write \(\Psi(\mathcal A)\) for that pattern's cap mass.  Consequently
\[
\begin{aligned}
 J_N(t)&=f_N(h(t)),\\
 f_N(u)&=u^N+
   \sum_{\mathcal A}
     C_0^{\,r}(1-u)^r u^{s_1+\cdots+s_r}\Psi(\mathcal A).
\end{aligned}
\]
Each actual labelled arrangement is counted once, with no additional
marked starting site.
The lengths satisfy \(\sum_i(s_i+g_i)=N\), so
\(r+\sum_i s_i\le N\).  Thus \(f_N\) is a polynomial of degree at
most \(N\) and \(f_N(1)=1\).  At first derivative order at \(u=1\),
only the full ring and the one-interval terms survive.  For a gap of
length \(g\in\{1,\ldots,N-1\}\), the one interval has \(N\) labelled
starting positions, its cap mass is \(K_N(g)\), and the derivative of
\((1-u)u^{N-g}\) at \(1\) is \(-1\).  We have proved the exact
normalization
\begin{equation}\label{eq:cyl-polynomial}
 \deg f_N\le N,\qquad
 f_N(1)=1,\qquad
 f_N'(1)=N-C_0N\sum_{g=1}^{N-1}K_N(g)=ND_N.
\end{equation}

\subsection{The strip operator}

Use the strip coordinate \(s=2t-3l\), and write
\[
 \alpha_{\rm s}:=\frac{\pi}{6l}=\frac43,\qquad
 \beta_{\rm s}:=\alpha_{\rm s}^{-1}=\frac34,\qquad
 S_a:=\{s\in\mathbb C:|\Re s|<a\}.
\]
The local straight weight in this coordinate is
\[
 h(t)=\frac{\cos s-\cos3l}{\cos s+\cos l}
 \quad\text{when }s=2t-3l.
\]
Define
\[
 j(s):=f_N\left(\frac{\cos s-\cos3l}{\cos s+\cos l}\right),
 \qquad
 G(s):=j(s-2l)+j(s+2l)-j(s).
\]
Then \(j(2t-3l)=J_N(t)\), and
\(G(2t-l)\) is the left side of
\eqref{eq:cyl-differentiated-fusion}.  Both \(j\) and \(G\) are even,
so \(G\) has zeros of order at least \(N\) at \(s=\pm l\).
Also \(j(i\infty)=G(i\infty)=1\).

The following representation depends only on the polynomial degree,
not on fusion.  This distinction will allow the same operator to be
used for higher winding derivatives.  For each fixed \(N\) and any
polynomial \(f\) of degree at most \(N\), put
\[
 j_f(s)=f\left(\frac{\cos s-\cos3l}{\cos s+\cos l}\right),
 \qquad
 G_f(s)=j_f(s-2l)+j_f(s+2l)-j_f(s).
\]
Define the even map and the convolution
\[
 B(s):=\frac{\cos(\alpha_{\rm s}s)-d_0}
             {\cos(\alpha_{\rm s}s)+d_0},
 \qquad d_0:=\cos(\alpha_{\rm s}l)=\frac{\sqrt3}{2},
 \qquad
 (\mathcal C_g F)(s):=\int_{\mathbb R}g(y)F(s+iy)\,dy,
\]
\[
 g(y):=\frac1{2\pi}
 \left(\frac{2^{-1/2}}{\cosh y-2^{-1/2}}-\frac1{\cosh y}\right).
\]
The map \(B\) identifies the quotient of \(S_{3l}\) by \(s\mapsto-s\)
conformally with the disk.  It sends \(\pm l\) to \(0\), and
\(B'(l)\ne0\).  The kernel \(g\) is real and even.  With
\(\widehat g(k)=\int g(y)e^{iky}\,dy\), residues over the \(2\pi i\)
period give
\[
 1-\widehat g(k)
 =1-\frac{\sinh(6lk)-\sinh(4lk)}{\sinh(8lk)}
 =(2\cosh(2lk)-1)\frac{\sinh(6lk)}{\sinh(8lk)}.
\]
The expressions at \(k=0\) are interpreted by continuity.
Furthermore,
\[
 \int_{\mathbb R}g=\frac14,\qquad p_g:=\int_{\mathbb R}|g|<\frac12.
\]
For the strict inequality, the negative part is zero on \(|y|\le1\)
and is bounded by
\((1-2^{-1/2})/(2\pi\cosh y)\) on \(|y|>1\); integrating this bound
and adding twice the negative mass to \(1/4\) gives \(p_g<1/2\).

With these definitions, the degree bound gives the representation
\begin{equation}\label{eq:cyl-strip-representation}
\begin{gathered}
 G_f(s)=(I-\mathcal C_g)q(B(s))+\lambda
 \quad(s\in S_{3l}),\\
 q(z)=\sum_{m=1}^N q_mz^m,
\end{gathered}
\end{equation}
for some \(q\) and constant \(\lambda\).

To prove \eqref{eq:cyl-strip-representation}, first Fourier transform
\(j_f(iy)-f(1)\), which decays exponentially at both ends.  Shifting
the contour through one real period \(16l\) crosses the mirrored poles
at \(7l,9l\), each of order at most \(N\).  Their residues give, on
\(S_{7l}\),
\[
 j_f(s)=f(1)+\int_{\mathbb R}
       \frac{H(k)}{\sinh(8lk)}e^{ks}\,dk,
 \qquad
 H(k)=P_e(k)\sinh(lk)+P_o(k)\cosh(lk),
\]
where \(P_e\) is even, \(P_o\) is odd, and both have degree at most
\(N-1\).  The displayed symbol therefore gives
\[
 G_f'=(I-\mathcal C_g)F,\qquad
 F(s):=\int_{\mathbb R}\frac{kH(k)}{\sinh(6lk)}e^{ks}\,dk.
\]

Here \(F\) lies in the holomorphic odd Bergman space on \(S_{3l}\),
with area norm \(\int_{S_{3l}}|F|^2\,dA\).  This space explains the
degree bound in \eqref{eq:cyl-strip-representation}.  Plancherel on
the vertical lines \(s=r+iy\), followed by integration over
\(-3l<r<3l\), gives the Fourier-side norm multiplier
\(\sinh(6lk)/k\), up to the Fourier normalization.  The line
transforms satisfy
\(\partial_r\widehat F=k\widehat F\) on interior lines in the sense
of distributions, by the Cauchy--Riemann equation.  In this norm the
span of the representers for jets of orders \(0,\ldots,N-1\) at \(l\),
restricted to odd functions, consists exactly of the preceding
numerators \(kH(k)\): even jet orders give the \(\sinh(lk)\) terms
and odd orders the \(\cosh(lk)\) terms.

Equivalently, every odd holomorphic function on the strip has an even
primitive, which is a function of \(B\).  For a primitive
\(p(B(s))\) with \(p(z)=\sum_{n\ge1}p_nz^n\), conformal change of
variables gives
\[
 \int_{S_{3l}}|(p(B))'|^2\,dA
       =2\pi\sum_{n\ge1}n|p_n|^2.
\]
This parametrizes the whole odd Bergman space.  Since \(B'(l)\ne0\),
the jets of \(F=(p(B))'\) of orders \(0,\ldots,N-1\) at \(l\) are
triangularly equivalent to the first \(N\) coefficients of \(p\).
Their representer span is therefore exactly the set of \((q(B))'\)
with \(\deg q\le N\) and
\(q(0)=0\).  Thus \(F=(q(B))'\), and integration in \(s\) proves
\eqref{eq:cyl-strip-representation}.

For later use, express the convolution on these coefficients.
Because \(g\) is even, \(\mathcal C_g B^m\) is even and hence a
holomorphic function of \(z=B(s)\).  With the local inverse branch
\(s(0)=l\), set
\[
\begin{aligned}
 \Phi_m(z)&:=(\mathcal C_g B^m)(s(z)),&
 M_{nm}&:=[z^n]\Phi_m(z)\quad(n,m\ge1),\\
 M_{0m}&:=\Phi_m(0)=(\mathcal C_g B^m)(l).&&
\end{aligned}
\]
Since \(|B(l+iy)|<1\), the defining convolution also gives
\(|M_{0m}|\le p_g\).
The matrix \(M\) is real and self-adjoint on
\[
 \ell^2(n):=\left\{v:\sum_{n\ge1}n|v_n|^2<\infty\right\},
 \qquad \|M\|_{\ell^2(n)\to\ell^2(n)}\le p_g.
\]
Indeed \((\mathcal C_g B^m)'=\mathcal C_g(B^m)'\), so differentiation
identifies its action with \(\mathcal C_g\) on the odd Bergman space.
Imaginary translations are unitary on the
full strip Bergman space; their average against the real even kernel
is self-adjoint, has norm at most \(\|g\|_1\), and preserves oddness.
In particular,
\[
 nM_{nm}=mM_{mn}.
\]
This Hilbert-space contraction does not assert contraction on
\(\ell^1\).

Apply \eqref{eq:cyl-strip-representation} to the first-derivative
function \(G\).  Its order-\(N\) zero at \(B=0\) gives
\[
 q_n=\sum_{m=1}^N M_{nm}q_m\quad(1\le n<N),
 \qquad
 \lambda=\sum_{m=1}^N M_{0m}q_m.
\]
The compressed Hilbert contraction makes \(I-(M_{nm})_{n,m<N}\)
invertible.  Define its unique normalized solution by
\[
 Q_N(z):=\sum_{m=1}^N c_mz^m,\qquad
 c_N=1,\qquad c_n=\sum_{m=1}^N M_{nm}c_m\quad(1\le n<N).
\]
If \(q_N=0\), the contraction would force \(q=0\), and the zero at
\(l\) would force \(G=0\), contrary to \(G(i\infty)=1\).  Thus \(q\)
is a nonzero multiple of \(Q_N\).  For a polynomial with zero
constant coefficient define
\[
\begin{aligned}
 (\mathcal Tq)(s)&:=(I-\mathcal C_g)q(B(s))
                      +(\mathcal C_g q(B))(l),\\
 \mathcal E(q)&:=(\mathcal Tq)(i\infty)
       =\frac34q(1)+\sum_mM_{0m}q_m.
\end{aligned}
\]
Writing \(\widetilde G:=\mathcal TQ_N\) and
\(E_N:=\mathcal E(Q_N)\), we obtain
\begin{equation}\label{eq:cyl-normalized-G}
 E_N\ne0,\qquad G(s)=\frac{\widetilde G(s)}{E_N}.
\end{equation}
The nonvanishing here holds for each fixed \(N\), by the existing
function \(G\).  We next need bounds on this normalizer as \(N\)
varies.

\subsection{Uniform control of the normalizer}

The Hilbert contraction has determined \(Q_N\), but a two-sided
estimate for \(G\) requires a strict bound on its lower
coefficients.  The following estimates for \(M\) will give that bound.

\paragraph{Matrix entries.}
For one constant \(C\), independent of \(n,m\ge1\),
\begin{equation}\label{eq:cyl-M-bounds}
\begin{gathered}
 |M_{nm}|\le\frac Cn
   \min\left\{\left(\frac nm\right)^{\beta_{\rm s}},
              \left(\frac mn\right)^{\beta_{\rm s}}\right\},
 \\
 |M_{n+1,m}-M_{nm}|\le\frac C{n^2},
 \qquad
 |M_{0m}|\le C m^{-\beta_{\rm s}}.
\end{gathered}
\end{equation}
Here is a contour proof.  Choose a fixed \(0<\delta<l/2\),
independently of \(n,m\), and continue
\(\mathcal C_g B^m\) from \(S_{3l}\) to \(S_{3l+\delta}\) by moving
the \(y\)-contour by at most \(2\delta i\), so that the argument of
\(B\) stays in \(|\Re s|\le3l-\delta\).  The kernel is holomorphic
for \(|\Im y|<2l\), so no pole of \(g\) is crossed.
On these contours,
\[
 |g(y)|\le C e^{-|\Re y|},\qquad
 |B(s)|\le \exp(-c_1e^{-\alpha_{\rm s}|\Im s|})
 \quad (|\Re s|\le3l-\delta).
\]
As an even function of \(s\), the continuation descends through the
inverse branch point \(B(0)\) to a
\(z\)-domain containing, for some \(\epsilon>0\),
\[
 |z|<1+\epsilon,\qquad
 |\arg(1-z)|<\frac{\pi}{2}+\epsilon,\qquad z\ne1.
\]
Away from \(z=1\) this follows by compactness and the disk mapping.
Near \(1\), inversion of
\(\cos(\alpha_{\rm s}s)=d_0(1+z)/(1-z)\) gives both the continued
domain and
\(e^{-\alpha_{\rm s}|\Im s|}\asymp |1-z|\).  Consequently,
\[
 |\Phi_m(z)|
 \le C\int_{\mathbb R}e^{-|y|}
       \exp\!\left(-c_2m|1-z|e^{-\alpha_{\rm s}|y|}\right)\,dy
 \le C'\min\{1,(m|1-z|)^{-\beta_{\rm s}}\}.
\]

For sufficiently large \(n\), use Cauchy's coefficient contour with
outer radius \(1+\epsilon/2\), indented along
\(1-z=r e^{\pm i(\pi/2+\epsilon/2)}\) to \(r=1/n\), then along the
small circle of that radius.  On the rays,
\(|z|^{-n-1}\le C e^{-c n r}\), since their angle is greater than
\(\pi/2\); the small circle has length \(O(1/n)\) and bounded
\(|z|^{-n-1}\).  The preceding bound on \(\Phi_m\) therefore gives
\[
 |M_{nm}|\le\frac Cn\min\{1,(n/m)^{\beta_{\rm s}}\}.
\]
Applying the same contour to \((1-z)\Phi_m(z)\) gives the difference
bound in \eqref{eq:cyl-M-bounds}.  For the finitely many smaller
\(n\), use a fixed circle \(|z|=r_*<1\); there \(|1-z|\) is bounded
away from zero, so the bound on \(\Phi_m\) is uniform in \(m\) and
gives the same estimates after increasing \(C\).  The symmetry
\(nM_{nm}=mM_{mn}\) supplies the additional
\((m/n)^{\beta_{\rm s}}\) factor, and the bound at \(z=0\) gives the
stated estimate for \(M_{0m}\).

\paragraph{The mass of a high column.}
The entry and difference bounds also imply the sharper limit
\begin{equation}\label{eq:cyl-column-limit}
 \sum_{n\ge1}|M_{nm}|\longrightarrow p_g
 \qquad(m\to\infty).
\end{equation}
To prove it, consider the finite signed measures
\[
 \mu_m:=\sum_{n\ge1}M_{nm}\,\delta_{n/m}
 \quad\text{on }(0,\infty).
\]
The entry bound makes their total variations uniformly bounded and
tight in compact ratio ranges.  More precisely, the total variation
on \(n/m<\eta\) is \(O(\eta^{\beta_{\rm s}})\), and that on
\(n/m>R\) is \(O(R^{-\beta_{\rm s}})\), uniformly in \(m\).

For fixed \(a>0\) and all sufficiently large \(m\), choose
\(Y_m\to\infty\) so that
\(B(iY_m)=e^{-a/m}\).  For each fixed \(y\),
\[
 B(i(Y_m+y))^m\longrightarrow
       \exp(-a e^{-\alpha_{\rm s}y}).
\]
Since \(0<B(iu)<1\) for real \(u\), dominated convergence and
\(M_{0m}\to0\) show
\[
 \int e^{-au}\,d\mu_m(u)
 =\Phi_m(e^{-a/m})-M_{0m}
 \longrightarrow
 \int_{\mathbb R}g(y)\exp(-a e^{-\alpha_{\rm s}y})\,dy.
\]
This is the Laplace transform of the smooth pushforward of
\(g(y)\,dy\) by \(u=e^{-\alpha_{\rm s}y}\), whose density is
\[
 \rho(u)=\frac1{\alpha_{\rm s}u}
         g\!\left(-\frac{\log u}{\alpha_{\rm s}}\right).
\]
Tightness and uniqueness of Laplace transforms give weak convergence
to \(\rho(u)\,du\).  To pass from signed convergence to total
variation, interpolate the values \(mM_{nm}\) linearly at \(u=n/m\).
On each compact interval in \((0,\infty)\), the entry bound makes
these interpolants bounded and the difference bound makes them
equicontinuous.  Every locally uniform subsequential limit must be
\(\rho\): for a compactly supported continuous test function
\(\varphi\), the identity
\[
 \int\varphi\,d\mu_m
 =\frac1m\sum_n\varphi(n/m)\,[mM_{nm}]
\]
is a Riemann sum for the interpolant limit, and its weak limit is
\(\int\varphi(u)\rho(u)\,du\).  Thus the convergence is locally uniform.
The tight total-variation tails then give
\[
 \sum_{n\ge1}|M_{nm}|\longrightarrow
 \int_0^\infty|\rho(u)|\,du=\int_{\mathbb R}|g(y)|\,dy=p_g,
\]
which proves \eqref{eq:cyl-column-limit}.

\paragraph{The lower coefficients of \(Q_N\).}
Write \(c=c^{(N)}\) as a sequence by setting \(c_m=0\) for \(m>N\).
We claim
\begin{equation}\label{eq:cyl-c-margin}
\begin{gathered}
 \sup_N\|c^{(N)}\|_1<\infty,\qquad
 c_m^{(N)}\longrightarrow0\ \text{for every fixed }m,\\
 \limsup_{N\to\infty}\sum_{m<N}|c_m|
       \le\frac{p_g}{1-p_g}<1.
\end{gathered}
\end{equation}
The proof uses the Hilbert contraction for possible fixed mass and
\eqref{eq:cyl-column-limit} for mass escaping to high indices.

Suppose first that \(\|c^{(N)}\|_1\) is unbounded along a subsequence
with \(N\to\infty\).
Divide by that norm to obtain \(\|v^{(N)}\|_1=1\).  The recurrence
for \(c\), including \(c_N=1\), becomes
\[
 \left\|v^{(N)}-\mathbf1_{\{n<N\}}Mv^{(N)}\right\|_1\longrightarrow0.
\]
Pass to a coordinatewise limit \(v\in\ell^1\), using Fatou's lemma.
Every fixed row
\(M_{nm}\) tends to zero as \(m\to\infty\), by
\eqref{eq:cyl-M-bounds}; this lets the recurrence pass to the limit,
giving \(v=Mv\).  Also \(|v_n|\le C\|v\|_1/n\), so
\[
 \sum_n n|v_n|^2\le C\|v\|_1\sum_n|v_n|<\infty.
\]
The strict contraction on \(\ell^2(n)\) forces \(v=0\).

Let \(d_m:=\sum_n|M_{nm}|\).  The entry bound makes \(d_m\) uniformly
bounded, and \eqref{eq:cyl-column-limit} gives
\(d_m\le p_g+\varepsilon<1\) for all sufficiently large \(m\).
The approximate recurrence and \(v_m^{(N)}\to0\) at each fixed index
would then imply
\[
 1\le \sum_{m\le J}d_m|v_m^{(N)}|
       +(p_g+\varepsilon)\|v^{(N)}\|_1+o(1)
 \longrightarrow p_g+\varepsilon<1
\]
for a suitable fixed \(J\), a contradiction.  This proves the first
claim of \eqref{eq:cyl-c-margin}.  With bounded norms, the same
coordinatewise-limit argument applies directly to \(c^{(N)}\);
every limit is an \(\ell^2(n)\) fixed point of \(M\), hence zero.
This proves the second claim.

Finally the recurrence and the column norms give
\[
 \sum_{n<N}|c_n|\le\sum_{n\ge1}|(Mc)_n|
                 \le\sum_m d_m|c_m|.
\]
The contribution of any fixed initial set of columns tends to zero.
Given \(\varepsilon>0\), choose the remaining columns so that
\(d_m\le p_g+\varepsilon\), let \(N\to\infty\), and then let
\(\varepsilon\downarrow0\).  With
\(\|c\|_1=1+\sum_{m<N}|c_m|\), this yields
\[
 \limsup_N\sum_{m<N}|c_m|
 \le p_g\left(1+\limsup_N\sum_{m<N}|c_m|\right),
\]
which is the third claim of \eqref{eq:cyl-c-margin}.  The same split
also gives
\[
 \limsup_N\|Mc\|_1\le\frac{p_g}{1-p_g}<1.
\]

All coefficients are real.  Hence \eqref{eq:cyl-c-margin} makes
\(Q_N(1)=1+\sum_{m<N}c_m\) positive and bounded away from zero for
all sufficiently large \(N\).  Moreover
\(\sum_mM_{0m}c_m\to0\): fixed coefficients \(c_m\) vanish, the
\(\ell^1\) norms are bounded, and \(M_{0m}\to0\).  The normalizer
therefore satisfies
\[
\begin{gathered}
 E_N=\frac34Q_N(1)+o(1),\qquad \sup_N|E_N|<\infty,\\
 \liminf_{N\to\infty}E_N
 \ge\frac34\,\frac{1-2p_g}{1-p_g}>0.
\end{gathered}
\]
Together with the fixed-\(N\) nonvanishing in
\eqref{eq:cyl-normalized-G}, this also gives a uniform positive lower
bound on \(|E_N|\).  The positivity of \(E_N\) itself has so far been
proved only for sufficiently large \(N\), which is all the
asymptotic argument needs.

The coefficient recurrence cancels every power below \(N\):
\[
 \widetilde G(s)
   =B(s)^N-\sum_{n\ge N}(Mc)_n B(s)^n.
\]
For real \(y\), \(z=B(iy)\) lies in \((0,1)\).  The last strict bound
on \(\|Mc\|_1\) and the positive bounds on \(E_N\) therefore give
\begin{equation}\label{eq:cyl-positive-G}
 G(iy)\asymp B(iy)^N>0
 \qquad(y\in\mathbb R,\ N\ge N_0),
\end{equation}
for some fixed \(N_0\), with constants uniform in \(y,N\).
This is the two-sided positive estimate needed for both masses.

\subsection{Recovering the two positive masses}

The two conclusions of Proposition~\ref{prop:cylinder} come from
different tails of the strip calculation.  The inverse shift kernel
at \(s=-l\) gives \(P_N\), while the slower tail of \(g\) determines
\(f_N'(1)=ND_N\).

The inverse of the shift symbol is
\[
 (2\cosh(2lk)-1)^{-1}=\frac{\cosh(lk)}{\cosh(3lk)}.
\]
Its Fourier kernel can be written
\[
 \mathcal K(y)=
 \frac{\operatorname{sech}(\alpha_{\rm s}y+i\pi/6)
       +\operatorname{sech}(\alpha_{\rm s}y-i\pi/6)}{12l}.
\]
It is analytic for \(|\Im y|<2l\).  Uniformly for \(v\) in any closed
inner subinterval of \((-2l,2l)\), its tails are
\[
\begin{aligned}
 \mathcal K(u+iv)&\sim c_h e^{-\alpha_{\rm s}u}e^{-i\alpha_{\rm s}v}
       &&(u\to+\infty),\\
 \mathcal K(u+iv)&\sim c_h e^{\alpha_{\rm s}u}e^{i\alpha_{\rm s}v}
       &&(u\to-\infty).
\end{aligned}
\qquad c_h=\frac{\sqrt3}{6l}>0.
\]
The endpoint behavior
\(j(s)-1=O_N(e^{-|\Im s|})\) on closed inner substrips of \(S_{7l}\)
justifies Fourier inversion of \(j-1\) and \(G-1\).  The constant
mode has multiplier \(1\), and \(\int\mathcal K=1\).  Contour
displacement to the imaginary axis, followed by taking real parts,
therefore gives
\[
 P_N=j(-l)=\int_{\mathbb R}
       \operatorname{Re}\mathcal K(y-il)\,G(iy)\,dy.
\]
Here \(G(iy)\) is real, and evenness permits the opposite sign of
the displacement as well.
Since \(\alpha_{\rm s}l=\pi/6\), the real part of the displaced
kernel is positively comparable to \(e^{-\alpha_{\rm s}|y|}\) for
large \(|y|\).  On the other hand
\[
 B(iy)=\frac{\cosh(\alpha_{\rm s}y)-d_0}
             {\cosh(\alpha_{\rm s}y)+d_0}
       =1-\Theta(e^{-\alpha_{\rm s}|y|})\in(0,1).
\]
On any bounded \(y\)-interval, \eqref{eq:cyl-positive-G} is
exponentially small in \(N\).  On either tail the same equation,
followed by the substitution \(v=e^{-\alpha_{\rm s}|y|}\), gives
\[
 \int_{|y|>Y}e^{-\alpha_{\rm s}|y|}B(iy)^N\,dy
 \asymp \int_0^{e^{-\alpha_{\rm s}Y}}e^{-\Theta(Nv)}\,dv
 \asymp N^{-1}
\]
for fixed sufficiently large \(Y\).  The positive tails dominate the
bounded middle, proving \(P_N\asymp N^{-1}\).

For the deficit, define
\[
 \kappa:=2(\cos l+\cos3l)(2\cos2l-1)>0,\qquad
 \gamma_g:=\frac{1-2^{-1/2}}{\pi}>0.
\]
At fixed \(N\), expansion of the rational expression for \(j\) at
imaginary infinity gives
\[
 \lim_{y\to\infty}e^y(1-G(iy))=\kappa f_N'(1).
\]
For \(Q_N\), put
\[
 D_Q(u):=Q_N(1)-Q_N(B(iu)),\qquad
 I_m:=\int_{\mathbb R}e^u(1-B(iu)^m)\,du.
\]
The constant terms cancel in the endpoint difference of
\eqref{eq:cyl-normalized-G}, leaving
\[
 1-G(iy)=\frac1{E_N}
 \left(D_Q(y)-\int_{\mathbb R}g(t)D_Q(y+t)\,dt\right).
\]
For fixed \(N\), \(D_Q(u)=O_N(e^{-\alpha_{\rm s}|u|})\).
Also \(g(v)\sim-\gamma_g e^{-|v|}\) and
\(e^y|g(u-y)|\le Ce^u\).  Since \(\alpha_{\rm s}>1\), the direct
term tends to zero after multiplication by \(e^y\), and dominated
convergence in the convolution gives the exact fixed-\(N\) limit
\[
 \kappa f_N'(1)=\frac{\gamma_g}{E_N}
                    \sum_{m=1}^N c_m I_m.
\]
No limit in \(N\) has been interchanged with this endpoint limit.

The integrals \(I_m\) are nonnegative and increase with \(m\), because
\(0<B(iu)<1\).  From
\(1-B(iu)\asymp e^{-\alpha_{\rm s}|u|}\) at infinity,
\[
 I_m\asymp m^{\beta_{\rm s}}.
\]
For example, on \(u\ge0\) the integrand is comparable to
\(e^u\min\{1,m e^{-\alpha_{\rm s}u}\}\); splitting at
\(u=\beta_{\rm s}\log m\) gives the stated power, and the negative
half-line is bounded.  Since \(c_N=1\),
\[
 \left(1-\sum_{m<N}|c_m|\right)I_N
 \le\sum_{m=1}^N c_m I_m
 \le\left(1+\sum_{m<N}|c_m|\right)I_N.
\]
The strict margin in \eqref{eq:cyl-c-margin} and the positive bounds
on \(E_N\) show
\[
 f_N'(1)\asymp N^{\beta_{\rm s}}=N^{3/4},
 \qquad
 D_N=\frac{f_N'(1)}N\asymp N^{-1/4}.
\]

These comparisons have been proved for sufficiently large \(N\).
For every remaining \(N\), the cap sum is finite and the full
horizontal polygon has positive weight \(h(l)^N\), so \(P_N>0\).
Also a half-plane path from \(0\) to \(N\) exists with positive weight.
Together with \eqref{eq:cyl-half-plane} and \(K_N(g)\le K_\infty(g)\),
this gives
\[
 D_N\ge 2\cos(3l)\sum_{g\ge N}K_\infty(g)>0.
\]
The finitely many remaining circumferences can
therefore be absorbed into the comparison constants, completing
Proposition~\ref{prop:cylinder}.

Finally the deficit also saturates the half-plane bound.  If
\(C_0=2\cos3l\), then
\[
 1-D_N=C_0\sum_{g<N}K_N(g)
 \le C_0\sum_{g<N}K_\infty(g)
 \le C_0\sum_{g>0}K_\infty(g)\le1.
\]
Since \(D_N\to0\), the last sum equals \(1\), and the same comparison
gives
\[
 C_0\sum_{g\ge N}K_\infty(g)
 =1-C_0\sum_{g<N}K_\infty(g)
 \le D_N=O(N^{-1/4}).
\]
Thus
\[
 2\cos(3l)\sum_{g>0}K_\infty(g)=1,\qquad
 \sum_{g\ge N}K_\infty(g)=O(N^{-1/4}),
\]
as required by the later boundary comparisons.
